\documentclass[11pt,a4paper]{article}

%=========================================================
% PACKAGES
%=========================================================
\usepackage[utf8]{inputenc}
\usepackage[T1]{fontenc}
\usepackage[french]{babel}
\usepackage[a4paper,margin=2.2cm]{geometry}
\usepackage{lmodern}
\usepackage{amsmath,amssymb,amsthm,mathtools}
\usepackage{fancyhdr}
\usepackage{lastpage}
\usepackage{enumitem}
\usepackage{array,tabularx}
\usepackage{tikz}
\usepackage[most]{tcolorbox}
\usepackage{hyperref}
\usepackage{xcolor}

%=========================================================
% MISE EN PAGE
%=========================================================
\setlength{\headheight}{14pt}
\pagestyle{fancy}
\fancyhf{}
\fancyhead[L]{Mathématiques CPGE}
\fancyhead[C]{Espaces euclidiens}
\fancyhead[R]{Page \thepage/\pageref{LastPage}}
\fancyfoot[C]{Cours complet -- Algèbre euclidienne}

\hypersetup{
    colorlinks=true,
    linkcolor=blue!50!black,
    urlcolor=blue!50!black
}

\setlist[itemize]{label=\(\triangleright\)}
\setlist[enumerate]{label=\textbf{\arabic*.}}

%=========================================================
% COULEURS
%=========================================================
\definecolor{bleufonce}{RGB}{25,65,130}
\definecolor{bleuclair}{RGB}{235,243,255}
\definecolor{vertfonce}{RGB}{30,120,70}
\definecolor{vertclair}{RGB}{235,250,240}
\definecolor{rougefonce}{RGB}{170,40,40}
\definecolor{rougeclair}{RGB}{255,238,238}
\definecolor{orangefonce}{RGB}{190,110,20}
\definecolor{orangeclair}{RGB}{255,246,230}
\definecolor{violetfonce}{RGB}{95,55,145}
\definecolor{violetclair}{RGB}{246,241,255}

\tcbset{
    enhanced,
    sharp corners,
    boxrule=0.7pt,
    before skip=8pt,
    after skip=8pt
}

\newtcolorbox{definitionbox}[1][]{
    colback=bleuclair,
    colframe=bleufonce,
    title=\textbf{Définition #1}
}

\newtcolorbox{theorembox}[1][]{
    colback=vertclair,
    colframe=vertfonce,
    title=\textbf{Théorème / Propriété #1}
}

\newtcolorbox{methodbox}[1][]{
    colback=orangeclair,
    colframe=orangefonce,
    title=\textbf{Méthode #1}
}

\newtcolorbox{retenirbox}[1][]{
    colback=violetclair,
    colframe=violetfonce,
    title=\textbf{À retenir #1}
}

\newtcolorbox{erreurbox}[1][]{
    colback=rougeclair,
    colframe=rougefonce,
    title=\textbf{Erreur classique #1}
}

\newtcolorbox{exercicebox}[1][]{
    colback=white,
    colframe=black!55,
    title=\textbf{Exercice #1}
}

\newtcolorbox{correctionbox}[1][]{
    colback=gray!5,
    colframe=black!65,
    title=\textbf{Correction #1}
}

%=========================================================
% COMMANDES
%=========================================================
\newcommand{\R}{\mathbb R}
\newcommand{\K}{\mathbb K}
\newcommand{\N}{\mathbb N}
\newcommand{\M}{\mathcal M}
\newcommand{\Lcal}{\mathcal L}
\newcommand{\Vect}{\operatorname{Vect}}
\newcommand{\Ker}{\operatorname{Ker}}
\newcommand{\Imf}{\operatorname{Im}}
\newcommand{\rg}{\operatorname{rg}}
\newcommand{\tr}{\operatorname{tr}}
\newcommand{\Id}{\operatorname{Id}}
\newcommand{\diag}{\operatorname{diag}}
\newcommand{\Mat}{\operatorname{Mat}}
\newcommand{\dist}{\operatorname{dist}}
\newcommand{\proj}{\operatorname{proj}}
\newcommand{\SO}{\operatorname{SO}}
\newcommand{\OO}{\operatorname{O}}
\newcommand{\norme}[1]{\left\lVert #1\right\rVert}
\newcommand{\abs}[1]{\left|#1\right|}
\newcommand{\scal}[2]{\left\langle #1,#2\right\rangle}

\newtheorem{prop}{Proposition}[section]
\newtheorem{lemme}[prop]{Lemme}
\newtheorem{corollaire}[prop]{Corollaire}

%=========================================================
% DOCUMENT
%=========================================================
\begin{document}

\begin{center}
    {\Huge \textbf{Espaces euclidiens}}\\[2mm]
    {\Large Produit scalaire, orthogonalité, projections, isométries et matrices symétriques}\\[2mm]
    {\large Cours complet pour classes préparatoires scientifiques}\\[2mm]
    \rule{0.88\linewidth}{0.5pt}
\end{center}

\vspace{0.4cm}

\tableofcontents

\newpage

%=========================================================
\section*{Introduction générale}
\addcontentsline{toc}{section}{Introduction générale}

Dans un espace vectoriel général, on sait parler de combinaisons linéaires, de bases, de dimension, d'applications linéaires et de réduction.  
Avec les espaces vectoriels normés, on sait aussi parler de longueur, de distance, de convergence et de continuité.

Mais en géométrie, une notion manque encore : l'angle.  
Pour parler rigoureusement d'orthogonalité, de projection, de distance à un sous-espace, de symétrie orthogonale, d'isométrie ou encore de matrice orthogonale, il faut enrichir l'espace vectoriel par un \textbf{produit scalaire}.

\begin{retenirbox}
Un espace euclidien est un espace vectoriel réel de dimension finie muni d'un produit scalaire.  
Le produit scalaire permet de faire de la géométrie dans un cadre algébrique.
\end{retenirbox}

Ce chapitre évite de refaire toute la théorie des espaces vectoriels normés.  
On utilisera seulement le fait qu'un produit scalaire définit une norme :
\[
\norme{x}=\sqrt{\scal{x}{x}},
\]
et donc une distance :
\[
d(x,y)=\norme{x-y}.
\]

%=========================================================
\section{Produit scalaire et espace euclidien}

\subsection{Produit scalaire}

\begin{definitionbox}[Produit scalaire réel]
Soit \(E\) un espace vectoriel réel.  
Un \textbf{produit scalaire} sur \(E\) est une application
\[
\scal{\cdot}{\cdot}:E\times E\longrightarrow \R
\]
vérifiant, pour tous \(x,y,z\in E\) et tout \(\lambda\in\R\) :

\begin{enumerate}
    \item \textbf{Bilinéarité :}
    \[
    \scal{\lambda x+y}{z}
    =
    \lambda\scal{x}{z}+\scal{y}{z},
    \]
    et
    \[
    \scal{x}{\lambda y+z}
    =
    \lambda\scal{x}{y}+\scal{x}{z}.
    \]

    \item \textbf{Symétrie :}
    \[
    \scal{x}{y}=\scal{y}{x}.
    \]

    \item \textbf{Positivité :}
    \[
    \scal{x}{x}\geq 0.
    \]

    \item \textbf{Définie positivité :}
    \[
    \scal{x}{x}=0\Longleftrightarrow x=0.
    \]
\end{enumerate}
\end{definitionbox}

\begin{definitionbox}[Espace euclidien]
Un \textbf{espace euclidien} est un espace vectoriel réel de dimension finie muni d'un produit scalaire.
\end{definitionbox}

\begin{erreurbox}
Un produit scalaire est toujours défini sur un espace vectoriel réel dans ce chapitre.  
Sur un espace vectoriel complexe, on parle plutôt de produit hermitien, avec une sesquilinéarité.
\end{erreurbox}

\subsection{Exemples fondamentaux}

\begin{theorembox}[Produit scalaire canonique de \(\R^n\)]
Sur \(\R^n\), l'application
\[
\scal{x}{y}=x_1y_1+\cdots+x_ny_n,
\]
où
\[
x=(x_1,\dots,x_n),\qquad y=(y_1,\dots,y_n),
\]
est un produit scalaire.
\end{theorembox}

\begin{proof}
La bilinéarité vient de la distributivité de la somme et du produit réel.  
La symétrie vient de la commutativité du produit réel :
\[
x_iy_i=y_ix_i.
\]
De plus :
\[
\scal{x}{x}=x_1^2+\cdots+x_n^2\geq0.
\]
Enfin :
\[
x_1^2+\cdots+x_n^2=0
\]
si et seulement si chaque terme positif est nul, donc si et seulement si :
\[
x_1=\cdots=x_n=0.
\]
Ainsi \(x=0\).  
L'application est donc un produit scalaire.
\end{proof}

\begin{theorembox}[Produit scalaire sur un espace de polynômes]
Sur \(\R_2[X]\), l'application
\[
\scal{P}{Q}=P(0)Q(0)+P(1)Q(1)+P(2)Q(2)
\]
est un produit scalaire.
\end{theorembox}

\begin{proof}
La bilinéarité et la symétrie sont immédiates.

Pour tout \(P\in\R_2[X]\),
\[
\scal{P}{P}=P(0)^2+P(1)^2+P(2)^2\geq0.
\]
Si :
\[
\scal{P}{P}=0,
\]
alors :
\[
P(0)=P(1)=P(2)=0.
\]
Le polynôme \(P\) est de degré au plus \(2\) et possède trois racines distinctes.  
Donc :
\[
P=0.
\]
L'application est donc définie positive.
\end{proof}

\begin{erreurbox}
Sur \(\R[X]\), la formule
\[
\scal{P}{Q}=P(0)Q(0)+P(1)Q(1)+P(2)Q(2)
\]
n'est pas un produit scalaire, car un polynôme non nul de degré supérieur ou égal à \(3\) peut s'annuler en \(0,1,2\).  
Par exemple :
\[
P(X)=X(X-1)(X-2).
\]
\end{erreurbox}

\subsection{Norme associée}

\begin{definitionbox}[Norme euclidienne]
Si \(E\) est un espace euclidien, la norme associée au produit scalaire est :
\[
\norme{x}=\sqrt{\scal{x}{x}}.
\]
\end{definitionbox}

\begin{retenirbox}
Dans un espace euclidien :
\[
\norme{x}^2=\scal{x}{x}.
\]
Cette égalité est utilisée constamment.
\end{retenirbox}

%=========================================================
\section{Inégalité de Cauchy-Schwarz et conséquences}

\subsection{Inégalité de Cauchy-Schwarz}

\begin{theorembox}[Cauchy-Schwarz]
Pour tous \(x,y\in E\),
\[
|\scal{x}{y}|\leq \norme{x}\norme{y}.
\]
De plus, il y a égalité si et seulement si \(x\) et \(y\) sont colinéaires.
\end{theorembox}

\begin{proof}
Si \(y=0\), l'inégalité est immédiate.

Supposons \(y\neq0\).  
Pour tout \(t\in\R\), on a :
\[
0\leq \norme{x-ty}^2.
\]
En développant :
\[
\norme{x-ty}^2
=
\scal{x-ty}{x-ty}
=
\norme{x}^2-2t\scal{x}{y}+t^2\norme{y}^2.
\]
Le polynôme réel
\[
P(t)=\norme{y}^2t^2-2\scal{x}{y}t+\norme{x}^2
\]
est positif pour tout réel \(t\).  
Son discriminant est donc négatif ou nul :
\[
\Delta
=
4\scal{x}{y}^2-4\norme{x}^2\norme{y}^2
\leq0.
\]
Ainsi :
\[
\scal{x}{y}^2\leq\norme{x}^2\norme{y}^2.
\]
En prenant la racine carrée :
\[
|\scal{x}{y}|\leq\norme{x}\norme{y}.
\]

Le cas d'égalité correspond à \(\Delta=0\), donc à l'existence d'un réel \(t_0\) tel que :
\[
\norme{x-t_0y}^2=0.
\]
Cela équivaut à :
\[
x=t_0y.
\]
Donc \(x\) et \(y\) sont colinéaires.  
La réciproque est immédiate.
\end{proof}

\subsection{Inégalité triangulaire}

\begin{theorembox}[Inégalité triangulaire]
Pour tous \(x,y\in E\),
\[
\norme{x+y}\leq\norme{x}+\norme{y}.
\]
\end{theorembox}

\begin{proof}
On développe :
\[
\norme{x+y}^2
=
\norme{x}^2+2\scal{x}{y}+\norme{y}^2.
\]
Par Cauchy-Schwarz :
\[
\scal{x}{y}\leq |\scal{x}{y}|\leq\norme{x}\norme{y}.
\]
Donc :
\[
\norme{x+y}^2
\leq
\norme{x}^2+2\norme{x}\norme{y}+\norme{y}^2
=
(\norme{x}+\norme{y})^2.
\]
Comme les deux membres sont positifs :
\[
\norme{x+y}\leq\norme{x}+\norme{y}.
\]
\end{proof}

\subsection{Identités remarquables}

\begin{theorembox}[Identité de polarisation]
Pour tous \(x,y\in E\),
\[
\scal{x}{y}
=
\frac12\left(\norme{x+y}^2-\norme{x}^2-\norme{y}^2\right).
\]
On a aussi :
\[
\scal{x}{y}
=
\frac14\left(\norme{x+y}^2-\norme{x-y}^2\right).
\]
\end{theorembox}

\begin{proof}
On développe :
\[
\norme{x+y}^2
=
\norme{x}^2+2\scal{x}{y}+\norme{y}^2.
\]
D'où la première formule.

De même :
\[
\norme{x-y}^2
=
\norme{x}^2-2\scal{x}{y}+\norme{y}^2.
\]
En soustrayant :
\[
\norme{x+y}^2-\norme{x-y}^2=4\scal{x}{y}.
\]
D'où la seconde formule.
\end{proof}

\begin{theorembox}[Identité du parallélogramme]
Pour tous \(x,y\in E\),
\[
\norme{x+y}^2+\norme{x-y}^2
=
2\norme{x}^2+2\norme{y}^2.
\]
\end{theorembox}

\begin{proof}
On additionne les deux développements :
\[
\norme{x+y}^2
=
\norme{x}^2+2\scal{x}{y}+\norme{y}^2,
\]
\[
\norme{x-y}^2
=
\norme{x}^2-2\scal{x}{y}+\norme{y}^2.
\]
Les termes en produit scalaire se compensent.
\end{proof}

\begin{methodbox}[Exploiter Cauchy-Schwarz]
Cauchy-Schwarz sert à :
\begin{enumerate}
    \item majorer un produit scalaire ;
    \item prouver qu'une expression est bien définie ;
    \item démontrer l'inégalité triangulaire ;
    \item reconnaître les cas d'égalité par colinéarité ;
    \item établir des inégalités sur des sommes ou des intégrales.
\end{enumerate}
\end{methodbox}

%=========================================================
\section{Orthogonalité}

\subsection{Vecteurs orthogonaux}

\begin{definitionbox}[Orthogonalité]
Deux vecteurs \(x,y\in E\) sont dits orthogonaux si :
\[
\scal{x}{y}=0.
\]
On note :
\[
x\perp y.
\]
\end{definitionbox}

\begin{theorembox}[Pythagore]
Si \(x\perp y\), alors :
\[
\norme{x+y}^2=\norme{x}^2+\norme{y}^2.
\]
\end{theorembox}

\begin{proof}
On développe :
\[
\norme{x+y}^2
=
\norme{x}^2+2\scal{x}{y}+\norme{y}^2.
\]
Si \(x\perp y\), alors \(\scal{x}{y}=0\), d'où :
\[
\norme{x+y}^2=\norme{x}^2+\norme{y}^2.
\]
\end{proof}

\begin{theorembox}[Pythagore généralisé]
Si \(x_1,\dots,x_p\) sont deux à deux orthogonaux, alors :
\[
\norme{x_1+\cdots+x_p}^2
=
\norme{x_1}^2+\cdots+\norme{x_p}^2.
\]
\end{theorembox}

\begin{proof}
On développe :
\[
\norme{x_1+\cdots+x_p}^2
=
\scal{\sum_{i=1}^p x_i}{\sum_{j=1}^p x_j}
=
\sum_{i=1}^p\sum_{j=1}^p\scal{x_i}{x_j}.
\]
Les termes avec \(i\neq j\) sont nuls par orthogonalité.  
Il reste :
\[
\sum_{i=1}^p\scal{x_i}{x_i}
=
\sum_{i=1}^p\norme{x_i}^2.
\]
\end{proof}

\subsection{Orthogonal d'une partie}

\begin{definitionbox}[Orthogonal d'une partie]
Soit \(A\subset E\).  
L'orthogonal de \(A\), noté \(A^\perp\), est :
\[
A^\perp=\{x\in E:\ \forall a\in A,\ \scal{x}{a}=0\}.
\]
\end{definitionbox}

\begin{theorembox}
Pour toute partie \(A\subset E\), \(A^\perp\) est un sous-espace vectoriel de \(E\).
\end{theorembox}

\begin{proof}
On a \(0\in A^\perp\), car :
\[
\scal{0}{a}=0
\]
pour tout \(a\in A\).

Soient \(x,y\in A^\perp\) et \(\lambda,\mu\in\R\).  
Pour tout \(a\in A\),
\[
\scal{\lambda x+\mu y}{a}
=
\lambda\scal{x}{a}+\mu\scal{y}{a}
=
0.
\]
Donc :
\[
\lambda x+\mu y\in A^\perp.
\]
Ainsi \(A^\perp\) est un sous-espace vectoriel.
\end{proof}

\begin{retenirbox}
Même si \(A\) n'est pas un sous-espace vectoriel, \(A^\perp\) est toujours un sous-espace vectoriel.
\end{retenirbox}

\subsection{Orthogonal d'une somme et inclusion}

\begin{theorembox}
Si \(A\subset B\), alors :
\[
B^\perp\subset A^\perp.
\]
\end{theorembox}

\begin{proof}
Si \(x\in B^\perp\), alors \(x\) est orthogonal à tout élément de \(B\).  
Comme \(A\subset B\), \(x\) est en particulier orthogonal à tout élément de \(A\).  
Donc :
\[
x\in A^\perp.
\]
\end{proof}

\begin{theorembox}
Si \(F\) et \(G\) sont deux sous-espaces vectoriels de \(E\), alors :
\[
(F+G)^\perp=F^\perp\cap G^\perp.
\]
\end{theorembox}

\begin{proof}
Soit \(x\in(F+G)^\perp\).  
Alors \(x\) est orthogonal à tout vecteur de \(F+G\).  
Comme \(F\subset F+G\) et \(G\subset F+G\), on a :
\[
x\in F^\perp
\quad\text{et}\quad
x\in G^\perp.
\]
Donc :
\[
x\in F^\perp\cap G^\perp.
\]

Réciproquement, soit \(x\in F^\perp\cap G^\perp\).  
Tout élément de \(F+G\) s'écrit \(f+g\), avec \(f\in F\) et \(g\in G\).  
Alors :
\[
\scal{x}{f+g}=\scal{x}{f}+\scal{x}{g}=0+0=0.
\]
Donc \(x\in(F+G)^\perp\).
\end{proof}

%=========================================================
\section{Familles orthogonales et bases orthonormées}

\subsection{Définitions}

\begin{definitionbox}[Famille orthogonale]
Une famille \((e_1,\dots,e_p)\) est orthogonale si :
\[
\forall i\neq j,\qquad \scal{e_i}{e_j}=0.
\]
\end{definitionbox}

\begin{definitionbox}[Famille orthonormée]
Une famille \((e_1,\dots,e_p)\) est orthonormée si :
\[
\forall i,j,\qquad \scal{e_i}{e_j}=\delta_{ij},
\]
c'est-à-dire :
\[
\scal{e_i}{e_j}=0\quad\text{si }i\neq j,
\]
et :
\[
\norme{e_i}=1\quad\text{pour tout }i.
\]
\end{definitionbox}

\subsection{Liberté}

\begin{theorembox}
Toute famille orthogonale de vecteurs non nuls est libre.
\end{theorembox}

\begin{proof}
Soit \((e_1,\dots,e_p)\) une famille orthogonale de vecteurs non nuls.  
Supposons :
\[
\lambda_1e_1+\cdots+\lambda_pe_p=0.
\]
On prend le produit scalaire avec \(e_k\).  
On obtient :
\[
\scal{\lambda_1e_1+\cdots+\lambda_pe_p}{e_k}=0.
\]
Par bilinéarité :
\[
\sum_{i=1}^p\lambda_i\scal{e_i}{e_k}=0.
\]
Par orthogonalité, tous les termes sont nuls sauf celui correspondant à \(i=k\).  
Donc :
\[
\lambda_k\norme{e_k}^2=0.
\]
Comme \(e_k\neq0\), on a \(\norme{e_k}^2>0\).  
Donc :
\[
\lambda_k=0.
\]
Ceci vaut pour tout \(k\), donc la famille est libre.
\end{proof}

\begin{corollaire}
Dans un espace euclidien de dimension \(n\), toute famille orthogonale de vecteurs non nuls possède au plus \(n\) vecteurs.
\end{corollaire}

\subsection{Coordonnées dans une base orthonormée}

\begin{theorembox}[Coordonnées dans une base orthonormée]
Soit \(\mathcal B=(e_1,\dots,e_n)\) une base orthonormée de \(E\).  
Alors, pour tout \(x\in E\),
\[
x=\sum_{k=1}^n \scal{x}{e_k}e_k.
\]
\end{theorembox}

\begin{proof}
Comme \(\mathcal B\) est une base, il existe des réels \(\lambda_1,\dots,\lambda_n\) tels que :
\[
x=\sum_{k=1}^n\lambda_ke_k.
\]
Prenons le produit scalaire avec \(e_j\).  
On obtient :
\[
\scal{x}{e_j}
=
\sum_{k=1}^n\lambda_k\scal{e_k}{e_j}.
\]
Comme la base est orthonormée :
\[
\scal{e_k}{e_j}=\delta_{kj}.
\]
Donc :
\[
\scal{x}{e_j}=\lambda_j.
\]
Ainsi :
\[
x=\sum_{k=1}^n\scal{x}{e_k}e_k.
\]
\end{proof}

\begin{corollaire}[Formule de Parseval]
Dans une base orthonormée \((e_1,\dots,e_n)\),
\[
\norme{x}^2
=
\sum_{k=1}^n\scal{x}{e_k}^2.
\]
\end{corollaire}

\begin{proof}
On écrit :
\[
x=\sum_{k=1}^n\scal{x}{e_k}e_k.
\]
La famille étant orthonormée, les vecteurs de cette somme sont deux à deux orthogonaux.  
Par Pythagore généralisé :
\[
\norme{x}^2
=
\sum_{k=1}^n\norme{\scal{x}{e_k}e_k}^2
=
\sum_{k=1}^n\scal{x}{e_k}^2\norme{e_k}^2.
\]
Comme \(\norme{e_k}=1\),
\[
\norme{x}^2
=
\sum_{k=1}^n\scal{x}{e_k}^2.
\]
\end{proof}

\begin{methodbox}[Utiliser une base orthonormée]
Dans une base orthonormée, les calculs deviennent simples :
\[
x=\sum_{k=1}^n \scal{x}{e_k}e_k,
\]
\[
\scal{x}{y}=\sum_{k=1}^n x_ky_k,
\]
\[
\norme{x}^2=\sum_{k=1}^n x_k^2.
\]
C'est pourquoi on cherche souvent à construire des bases orthonormées.
\end{methodbox}

%=========================================================
\section{Procédé de Gram-Schmidt}

\subsection{Énoncé}

\begin{theorembox}[Orthogonalisation de Gram-Schmidt]
Soit \((u_1,\dots,u_n)\) une famille libre de \(E\).  
On définit :
\[
v_1=u_1,
\]
puis, pour \(k\geq2\),
\[
v_k
=
u_k
-
\sum_{j=1}^{k-1}
\frac{\scal{u_k}{v_j}}{\norme{v_j}^2}v_j.
\]
Alors \((v_1,\dots,v_n)\) est une famille orthogonale de vecteurs non nuls et, pour tout \(k\),
\[
\Vect(v_1,\dots,v_k)=\Vect(u_1,\dots,u_k).
\]
En posant :
\[
e_k=\frac{v_k}{\norme{v_k}},
\]
on obtient une famille orthonormée.
\end{theorembox}

\begin{proof}
On raisonne par récurrence.

Pour \(k=1\), on a \(v_1=u_1\neq0\), et :
\[
\Vect(v_1)=\Vect(u_1).
\]

Supposons \(v_1,\dots,v_{k-1}\) construits, orthogonaux, non nuls, et tels que :
\[
\Vect(v_1,\dots,v_{k-1})=\Vect(u_1,\dots,u_{k-1}).
\]
On définit :
\[
v_k
=
u_k
-
\sum_{j=1}^{k-1}
\frac{\scal{u_k}{v_j}}{\norme{v_j}^2}v_j.
\]

Montrons que \(v_k\) est orthogonal à \(v_i\) pour \(i<k\).  
On calcule :
\[
\scal{v_k}{v_i}
=
\scal{u_k}{v_i}
-
\sum_{j=1}^{k-1}
\frac{\scal{u_k}{v_j}}{\norme{v_j}^2}
\scal{v_j}{v_i}.
\]
Comme les \(v_j\) sont orthogonaux, \(\scal{v_j}{v_i}=0\) si \(j\neq i\), et :
\[
\scal{v_i}{v_i}=\norme{v_i}^2.
\]
Donc :
\[
\scal{v_k}{v_i}
=
\scal{u_k}{v_i}
-
\frac{\scal{u_k}{v_i}}{\norme{v_i}^2}
\norme{v_i}^2
=0.
\]

Montrons que \(v_k\neq0\).  
Si \(v_k=0\), alors :
\[
u_k
=
\sum_{j=1}^{k-1}
\frac{\scal{u_k}{v_j}}{\norme{v_j}^2}v_j
\in
\Vect(v_1,\dots,v_{k-1})
=
\Vect(u_1,\dots,u_{k-1}),
\]
ce qui contredit la liberté de \((u_1,\dots,u_k)\).  
Donc \(v_k\neq0\).

Enfin, comme \(v_k\) est obtenu à partir de \(u_k\) en ajoutant une combinaison linéaire de \(v_1,\dots,v_{k-1}\), on a :
\[
\Vect(v_1,\dots,v_k)=\Vect(u_1,\dots,u_k).
\]
\end{proof}

\begin{corollaire}[Existence de bases orthonormées]
Tout espace euclidien admet une base orthonormée.
\end{corollaire}

\begin{proof}
On part d'une base quelconque de \(E\), puis on applique le procédé de Gram-Schmidt et on normalise les vecteurs obtenus.
\end{proof}

\begin{methodbox}[Appliquer Gram-Schmidt]
Pour orthonormaliser une famille libre \((u_1,\dots,u_n)\) :
\begin{enumerate}
    \item poser \(v_1=u_1\), puis \(e_1=\dfrac{v_1}{\norme{v_1}}\) ;
    \item retirer à \(u_2\) sa projection sur \(\Vect(v_1)\) ;
    \item retirer à \(u_k\) sa projection sur \(\Vect(v_1,\dots,v_{k-1})\) ;
    \item normaliser chaque \(v_k\).
\end{enumerate}
\end{methodbox}

\begin{erreurbox}
Dans Gram-Schmidt, on projette sur les vecteurs déjà orthogonaux \(v_j\), pas directement sur les vecteurs initiaux \(u_j\), sauf au premier rang.
\end{erreurbox}

%=========================================================
\section{Orthogonal d'un sous-espace et somme directe orthogonale}

\subsection{Théorème fondamental}

\begin{theorembox}[Décomposition orthogonale]
Soit \(F\) un sous-espace vectoriel d'un espace euclidien \(E\).  
Alors :
\[
E=F\oplus F^\perp.
\]
\end{theorembox}

\begin{proof}
Comme \(F\) est un espace euclidien pour le produit scalaire induit, il admet une base orthonormée :
\[
(e_1,\dots,e_p).
\]
On complète cette famille orthonormée en une base orthonormée de \(E\) :
\[
(e_1,\dots,e_p,e_{p+1},\dots,e_n).
\]
Alors :
\[
F=\Vect(e_1,\dots,e_p).
\]
Montrons que :
\[
F^\perp=\Vect(e_{p+1},\dots,e_n).
\]

D'abord, tout vecteur de \(\Vect(e_{p+1},\dots,e_n)\) est orthogonal à \(F\), car la base est orthonormée.  
Donc :
\[
\Vect(e_{p+1},\dots,e_n)\subset F^\perp.
\]

Réciproquement, soit \(x\in F^\perp\).  
Écrivons :
\[
x=\sum_{k=1}^n\lambda_ke_k.
\]
Pour \(1\leq i\leq p\), comme \(x\in F^\perp\),
\[
0=\scal{x}{e_i}=\lambda_i.
\]
Donc :
\[
x=\sum_{k=p+1}^n\lambda_ke_k.
\]
Ainsi :
\[
F^\perp=\Vect(e_{p+1},\dots,e_n).
\]

On obtient donc :
\[
E=F\oplus F^\perp.
\]
\end{proof}

\begin{corollaire}
Pour tout sous-espace \(F\) de \(E\),
\[
\dim(F)+\dim(F^\perp)=\dim(E).
\]
\end{corollaire}

\begin{corollaire}
Pour tout sous-espace \(F\) de \(E\),
\[
(F^\perp)^\perp=F.
\]
\end{corollaire}

\begin{proof}
On a toujours :
\[
F\subset (F^\perp)^\perp.
\]
En effet, tout vecteur de \(F\) est orthogonal à tout vecteur de \(F^\perp\).

De plus :
\[
\dim((F^\perp)^\perp)
=
\dim(E)-\dim(F^\perp).
\]
Or :
\[
\dim(F^\perp)=\dim(E)-\dim(F).
\]
Donc :
\[
\dim((F^\perp)^\perp)=\dim(F).
\]
L'inclusion entre deux sous-espaces de même dimension est donc une égalité.
\end{proof}

\begin{retenirbox}
Dans un espace euclidien de dimension finie :
\[
E=F\oplus F^\perp,
\qquad
\dim(F^\perp)=\dim(E)-\dim(F),
\qquad
(F^\perp)^\perp=F.
\]
\end{retenirbox}

%=========================================================
\section{Projection orthogonale}

\subsection{Définition}

\begin{definitionbox}[Projection orthogonale]
Soit \(F\) un sous-espace vectoriel de \(E\).  
Comme :
\[
E=F\oplus F^\perp,
\]
tout vecteur \(x\in E\) s'écrit de manière unique :
\[
x=p_F(x)+r_F(x),
\]
avec :
\[
p_F(x)\in F,
\qquad
r_F(x)\in F^\perp.
\]
Le vecteur \(p_F(x)\) est appelé le \textbf{projeté orthogonal} de \(x\) sur \(F\).
\end{definitionbox}

\begin{theorembox}[Caractérisation du projeté orthogonal]
Le vecteur \(p_F(x)\) est l'unique vecteur \(y\in F\) tel que :
\[
x-y\in F^\perp.
\]
\end{theorembox}

\begin{proof}
Par définition, dans la décomposition :
\[
x=p_F(x)+r_F(x),
\]
on a :
\[
p_F(x)\in F,
\qquad
r_F(x)=x-p_F(x)\in F^\perp.
\]
Donc \(p_F(x)\) vérifie bien la propriété.

Réciproquement, si \(y\in F\) vérifie :
\[
x-y\in F^\perp,
\]
alors :
\[
x=y+(x-y)
\]
est une décomposition de \(x\) comme somme d'un élément de \(F\) et d'un élément de \(F^\perp\).  
Par unicité de la décomposition dans la somme directe :
\[
E=F\oplus F^\perp,
\]
on a :
\[
y=p_F(x).
\]
\end{proof}

\subsection{Formule dans une base orthonormée}

\begin{theorembox}[Formule de projection]
Si \((e_1,\dots,e_p)\) est une base orthonormée de \(F\), alors :
\[
p_F(x)=\sum_{k=1}^p \scal{x}{e_k}e_k.
\]
\end{theorembox}

\begin{proof}
Posons :
\[
y=\sum_{k=1}^p \scal{x}{e_k}e_k.
\]
Alors \(y\in F\).  
Montrons que \(x-y\in F^\perp\).  
Comme \((e_1,\dots,e_p)\) est une base de \(F\), il suffit de montrer que \(x-y\) est orthogonal à chaque \(e_j\).

On calcule :
\[
\scal{x-y}{e_j}
=
\scal{x}{e_j}
-
\sum_{k=1}^p\scal{x}{e_k}\scal{e_k}{e_j}.
\]
Comme la base est orthonormée :
\[
\scal{e_k}{e_j}=\delta_{kj}.
\]
Donc :
\[
\scal{x-y}{e_j}
=
\scal{x}{e_j}-\scal{x}{e_j}=0.
\]
Ainsi \(x-y\in F^\perp\).  
Par caractérisation du projeté orthogonal :
\[
y=p_F(x).
\]
\end{proof}

\subsection{Projection sur une droite}

\begin{corollaire}
Si \(F=\R a\), avec \(a\neq0\), alors :
\[
p_F(x)=\frac{\scal{x}{a}}{\norme{a}^2}a.
\]
\end{corollaire}

\begin{proof}
Une base orthonormée de \(\R a\) est formée du vecteur :
\[
e=\frac{a}{\norme{a}}.
\]
Donc :
\[
p_F(x)=\scal{x}{e}e.
\]
Or :
\[
\scal{x}{e}=\frac{\scal{x}{a}}{\norme{a}},
\]
et :
\[
e=\frac{a}{\norme{a}}.
\]
Donc :
\[
p_F(x)=\frac{\scal{x}{a}}{\norme{a}^2}a.
\]
\end{proof}

\begin{methodbox}[Calculer un projeté orthogonal]
Pour calculer \(p_F(x)\) :
\begin{enumerate}
    \item trouver une base de \(F\) ;
    \item si elle n'est pas orthonormée, l'orthonormaliser avec Gram-Schmidt ;
    \item appliquer :
    \[
    p_F(x)=\sum_{k=1}^p\scal{x}{e_k}e_k.
    \]
\end{enumerate}
\end{methodbox}

\begin{methodbox}[Autre méthode : système normal]
Si \(F=\Vect(u_1,\dots,u_p)\), on peut chercher :
\[
p_F(x)=\lambda_1u_1+\cdots+\lambda_pu_p.
\]
La condition :
\[
x-p_F(x)\in F^\perp
\]
équivaut au système :
\[
\forall i,\qquad
\scal{x-\sum_{j=1}^p\lambda_ju_j}{u_i}=0.
\]
C'est un système linéaire en \(\lambda_1,\dots,\lambda_p\).
\end{methodbox}

%=========================================================
\section{Distance à un sous-espace}

\subsection{Définition}

\begin{definitionbox}[Distance à un sous-espace]
Soit \(F\) un sous-espace de \(E\).  
La distance de \(x\in E\) à \(F\) est :
\[
d(x,F)=\inf_{y\in F}\norme{x-y}.
\]
\end{definitionbox}

\subsection{Théorème du projeté orthogonal}

\begin{theorembox}[Distance à un sous-espace]
Pour tout \(x\in E\),
\[
d(x,F)=\norme{x-p_F(x)}.
\]
De plus, \(p_F(x)\) est l'unique point de \(F\) qui réalise cette distance.
\end{theorembox}

\begin{proof}
Soit \(y\in F\).  
On écrit :
\[
x-y=(x-p_F(x))+(p_F(x)-y).
\]
Or :
\[
x-p_F(x)\in F^\perp,
\]
et :
\[
p_F(x)-y\in F.
\]
Ces deux vecteurs sont donc orthogonaux.  
Par le théorème de Pythagore :
\[
\norme{x-y}^2
=
\norme{x-p_F(x)}^2+\norme{p_F(x)-y}^2.
\]
Donc :
\[
\norme{x-y}^2\geq \norme{x-p_F(x)}^2.
\]
Ainsi :
\[
\norme{x-y}\geq \norme{x-p_F(x)}.
\]
L'égalité a lieu si et seulement si :
\[
\norme{p_F(x)-y}^2=0,
\]
c'est-à-dire :
\[
y=p_F(x).
\]
Donc le minimum est atteint uniquement en \(p_F(x)\), et :
\[
d(x,F)=\norme{x-p_F(x)}.
\]
\end{proof}

\begin{retenirbox}
Le projeté orthogonal est le point du sous-espace le plus proche de \(x\).  
La distance à \(F\) est la norme du résidu orthogonal :
\[
x-p_F(x).
\]
\end{retenirbox}

%=========================================================
\section{Matrices de Gram et produit scalaire en base quelconque}

\subsection{Matrice de Gram}

\begin{definitionbox}[Matrice de Gram]
Soit \(\mathcal B=(e_1,\dots,e_n)\) une base de \(E\).  
La matrice de Gram de \(\mathcal B\) est :
\[
G_{\mathcal B}=(\scal{e_i}{e_j})_{1\leq i,j\leq n}.
\]
\end{definitionbox}

\begin{theorembox}
La matrice de Gram est symétrique définie positive.
\end{theorembox}

\begin{proof}
Elle est symétrique car :
\[
(G_{\mathcal B})_{ij}=\scal{e_i}{e_j}=\scal{e_j}{e_i}=(G_{\mathcal B})_{ji}.
\]

Soit \(X=(x_1,\dots,x_n)^T\neq0\).  
Posons :
\[
x=\sum_{i=1}^n x_ie_i.
\]
Comme \(X\neq0\) et \(\mathcal B\) est une base, on a \(x\neq0\).  
Alors :
\[
X^TG_{\mathcal B}X
=
\sum_{i,j}x_ix_j\scal{e_i}{e_j}
=
\scal{\sum_i x_ie_i}{\sum_j x_je_j}
=
\scal{x}{x}
=
\norme{x}^2>0.
\]
Donc \(G_{\mathcal B}\) est définie positive.
\end{proof}

\begin{theorembox}[Expression matricielle du produit scalaire]
Si \(X\) et \(Y\) sont les colonnes des coordonnées de \(x\) et \(y\) dans la base \(\mathcal B\), alors :
\[
\scal{x}{y}=X^TG_{\mathcal B}Y.
\]
\end{theorembox}

\begin{proof}
Écrivons :
\[
x=\sum_{i=1}^n x_ie_i,
\qquad
y=\sum_{j=1}^n y_je_j.
\]
Alors :
\[
\scal{x}{y}
=
\sum_{i=1}^n\sum_{j=1}^n x_iy_j\scal{e_i}{e_j}.
\]
C'est exactement le produit matriciel :
\[
X^TG_{\mathcal B}Y.
\]
\end{proof}

\begin{retenirbox}
Dans une base orthonormée :
\[
G_{\mathcal B}=I_n.
\]
Donc :
\[
\scal{x}{y}=X^TY.
\]
\end{retenirbox}

%=========================================================
\section{Formes linéaires et représentation par produit scalaire}

\begin{theorembox}[Représentation des formes linéaires]
Soit \(E\) un espace euclidien.  
Pour toute forme linéaire \(\varphi:E\to\R\), il existe un unique vecteur \(a\in E\) tel que :
\[
\forall x\in E,\qquad \varphi(x)=\scal{x}{a}.
\]
\end{theorembox}

\begin{proof}
Soit \((e_1,\dots,e_n)\) une base orthonormée de \(E\).  
Posons :
\[
a=\sum_{k=1}^n \varphi(e_k)e_k.
\]
Pour tout \(x\in E\), on écrit :
\[
x=\sum_{k=1}^n x_ke_k.
\]
Alors :
\[
\varphi(x)
=
\sum_{k=1}^n x_k\varphi(e_k).
\]
D'autre part :
\[
\scal{x}{a}
=
\scal{\sum_{k=1}^n x_ke_k}{\sum_{j=1}^n\varphi(e_j)e_j}
=
\sum_{k=1}^n x_k\varphi(e_k).
\]
Donc :
\[
\varphi(x)=\scal{x}{a}.
\]

Montrons l'unicité.  
Si \(a,b\in E\) conviennent, alors :
\[
\forall x\in E,\quad \scal{x}{a}=\scal{x}{b}.
\]
Donc :
\[
\forall x\in E,\quad \scal{x}{a-b}=0.
\]
En prenant \(x=a-b\), on obtient :
\[
\norme{a-b}^2=0.
\]
Donc :
\[
a=b.
\]
\end{proof}

\begin{corollaire}[Équation d'un hyperplan]
Tout hyperplan \(H\) d'un espace euclidien \(E\) peut s'écrire :
\[
H=a^\perp=\{x\in E:\scal{x}{a}=0\},
\]
pour un certain vecteur non nul \(a\).
\end{corollaire}

\begin{proof}
Un hyperplan est le noyau d'une forme linéaire non nulle \(\varphi\).  
D'après le théorème précédent, il existe \(a\neq0\) tel que :
\[
\varphi(x)=\scal{x}{a}.
\]
Donc :
\[
H=\Ker(\varphi)=\{x\in E:\scal{x}{a}=0\}=a^\perp.
\]
\end{proof}

%=========================================================
\section{Matrices orthogonales et isométries}

\subsection{Matrices orthogonales}

\begin{definitionbox}[Matrice orthogonale]
Une matrice \(A\in\M_n(\R)\) est orthogonale si :
\[
A^TA=I_n.
\]
On note l'ensemble des matrices orthogonales :
\[
O(n).
\]
\end{definitionbox}

\begin{theorembox}[Caractérisations]
Pour \(A\in\M_n(\R)\), les propriétés suivantes sont équivalentes :
\begin{enumerate}
    \item \(A\) est orthogonale ;
    \item \(A^TA=I_n\) ;
    \item \(AA^T=I_n\) ;
    \item les colonnes de \(A\) forment une base orthonormée de \(\R^n\) ;
    \item les lignes de \(A\) forment une base orthonormée de \(\R^n\).
\end{enumerate}
\end{theorembox}

\begin{proof}
La matrice \(A^TA\) a pour coefficient \((i,j)\) le produit scalaire de la colonne \(i\) de \(A\) avec la colonne \(j\) de \(A\).  
Ainsi :
\[
A^TA=I_n
\]
si et seulement si les colonnes de \(A\) forment une famille orthonormée.  
Comme il y en a \(n\) dans \(\R^n\), elles forment une base orthonormée.

Si \(A^TA=I_n\), alors \(A\) est inversible et :
\[
A^{-1}=A^T.
\]
Donc :
\[
AA^T=I_n.
\]
La caractérisation par les lignes s'obtient de la même façon à partir de \(AA^T=I_n\).
\end{proof}

\begin{corollaire}
Si \(A\in O(n)\), alors :
\[
\det(A)=\pm1.
\]
\end{corollaire}

\begin{proof}
Comme :
\[
A^TA=I_n,
\]
en prenant les déterminants :
\[
\det(A^T)\det(A)=1.
\]
Or :
\[
\det(A^T)=\det(A).
\]
Donc :
\[
\det(A)^2=1.
\]
Ainsi :
\[
\det(A)=\pm1.
\]
\end{proof}

\begin{definitionbox}[Groupe spécial orthogonal]
On note :
\[
SO(n)=\{A\in O(n):\det(A)=1\}.
\]
\end{definitionbox}

\subsection{Isométries vectorielles}

\begin{definitionbox}[Isométrie vectorielle]
Soit \(E\) un espace euclidien.  
Un endomorphisme \(u\in\Lcal(E)\) est une isométrie vectorielle si :
\[
\forall x\in E,\qquad \norme{u(x)}=\norme{x}.
\]
\end{definitionbox}

\begin{theorembox}[Caractérisations des isométries]
Pour \(u\in\Lcal(E)\), les propriétés suivantes sont équivalentes :
\begin{enumerate}
    \item \(u\) conserve la norme :
    \[
    \forall x\in E,\quad \norme{u(x)}=\norme{x};
    \]
    \item \(u\) conserve le produit scalaire :
    \[
    \forall x,y\in E,\quad \scal{u(x)}{u(y)}=\scal{x}{y};
    \]
    \item l'image d'une base orthonormée par \(u\) est une base orthonormée ;
    \item dans une base orthonormée, la matrice de \(u\) est orthogonale.
\end{enumerate}
\end{theorembox}

\begin{proof}
\textbf{1 implique 2.}  
On utilise l'identité de polarisation :
\[
\scal{u(x)}{u(y)}
=
\frac14\left(\norme{u(x)+u(y)}^2-\norme{u(x)-u(y)}^2\right).
\]
Comme \(u\) est linéaire :
\[
u(x)+u(y)=u(x+y),
\qquad
u(x)-u(y)=u(x-y).
\]
Comme \(u\) conserve la norme :
\[
\norme{u(x+y)}=\norme{x+y},
\qquad
\norme{u(x-y)}=\norme{x-y}.
\]
Donc :
\[
\scal{u(x)}{u(y)}
=
\frac14\left(\norme{x+y}^2-\norme{x-y}^2\right)
=
\scal{x}{y}.
\]

\textbf{2 implique 1.}  
En prenant \(y=x\), on obtient :
\[
\norme{u(x)}^2=\scal{u(x)}{u(x)}=\scal{x}{x}=\norme{x}^2.
\]
Donc :
\[
\norme{u(x)}=\norme{x}.
\]

\textbf{2 implique 3.}  
Si \((e_1,\dots,e_n)\) est orthonormée, alors :
\[
\scal{u(e_i)}{u(e_j)}=\scal{e_i}{e_j}=\delta_{ij}.
\]
Donc \((u(e_1),\dots,u(e_n))\) est orthonormée.  
Elle contient \(n\) vecteurs dans un espace de dimension \(n\), donc c'est une base.

\textbf{3 implique 4.}  
Dans une base orthonormée, les colonnes de la matrice de \(u\) sont les coordonnées des vecteurs \(u(e_j)\).  
Si ces vecteurs forment une base orthonormée, alors les colonnes de la matrice sont orthonormées.  
Donc la matrice est orthogonale.

\textbf{4 implique 2.}  
Si \(A\) est la matrice de \(u\) dans une base orthonormée, alors :
\[
A^TA=I_n.
\]
Pour des colonnes de coordonnées \(X,Y\), on a :
\[
\scal{u(x)}{u(y)}
=
(AX)^T(AY)
=
X^TA^TAY
=
X^TY
=
\scal{x}{y}.
\]
\end{proof}

\begin{retenirbox}
Une isométrie vectorielle conserve simultanément :
\[
\text{normes, distances, produits scalaires, angles et orthogonalité.}
\]
\end{retenirbox}

\subsection{Stabilité de l'orthogonal}

\begin{theorembox}
Soit \(u\) une isométrie vectorielle de \(E\).  
Si \(F\) est un sous-espace stable par \(u\), alors \(F^\perp\) est stable par \(u\).
\end{theorembox}

\begin{proof}
Soit \(x\in F^\perp\).  
On veut montrer que \(u(x)\in F^\perp\).

Soit \(y\in F\).  
Comme \(F\) est stable par \(u\), on a :
\[
u^{-1}(y)\in F,
\]
car une isométrie est bijective et \(F\) est stable par \(u\), donc stable par \(u^{-1}\).  
Comme \(x\in F^\perp\),
\[
\scal{x}{u^{-1}(y)}=0.
\]
Or \(u\) conserve le produit scalaire :
\[
\scal{u(x)}{y}
=
\scal{u(x)}{u(u^{-1}(y))}
=
\scal{x}{u^{-1}(y)}
=
0.
\]
Donc \(u(x)\) est orthogonal à tout \(y\in F\), d'où :
\[
u(x)\in F^\perp.
\]
\end{proof}

%=========================================================
\section{Projections et symétries orthogonales}

\subsection{Projecteurs orthogonaux}

\begin{definitionbox}[Projecteur orthogonal]
Un projecteur \(p\) est dit orthogonal si :
\[
\Ker(p)=\Imf(p)^\perp.
\]
De manière équivalente, \(p\) est la projection orthogonale sur \(\Imf(p)\).
\end{definitionbox}

\begin{theorembox}[Caractérisation]
Un endomorphisme \(p\) est un projecteur orthogonal si et seulement si :
\[
p^2=p
\]
et :
\[
\forall x,y\in E,\qquad \scal{p(x)}{y}=\scal{x}{p(y)}.
\]
\end{theorembox}

\begin{proof}
Supposons \(p\) projecteur orthogonal sur \(F\).  
Alors \(p^2=p\).  
Écrivons :
\[
x=p(x)+r_x,
\qquad r_x\in F^\perp,
\]
et :
\[
y=p(y)+r_y,
\qquad r_y\in F^\perp.
\]
Comme \(p(x),p(y)\in F\), et \(r_x,r_y\in F^\perp\), on a :
\[
\scal{p(x)}{y}
=
\scal{p(x)}{p(y)+r_y}
=
\scal{p(x)}{p(y)}.
\]
De même :
\[
\scal{x}{p(y)}
=
\scal{p(x)+r_x}{p(y)}
=
\scal{p(x)}{p(y)}.
\]
Donc :
\[
\scal{p(x)}{y}=\scal{x}{p(y)}.
\]

Réciproquement, supposons \(p^2=p\) et la relation de symétrie.  
On sait que :
\[
E=\Imf(p)\oplus\Ker(p).
\]
Montrons que :
\[
\Ker(p)\subset\Imf(p)^\perp.
\]
Soit \(x\in\Ker(p)\) et \(y\in\Imf(p)\).  
Il existe \(z\in E\) tel que :
\[
y=p(z).
\]
Alors :
\[
\scal{x}{y}
=
\scal{x}{p(z)}
=
\scal{p(x)}{z}
=
\scal{0}{z}
=
0.
\]
Donc :
\[
x\in\Imf(p)^\perp.
\]
Comme :
\[
\dim(\Ker(p))+\dim(\Imf(p))=\dim(E),
\]
et :
\[
\dim(\Imf(p)^\perp)=\dim(E)-\dim(\Imf(p)),
\]
les deux sous-espaces ont même dimension.  
Donc :
\[
\Ker(p)=\Imf(p)^\perp.
\]
\end{proof}

\subsection{Symétries orthogonales}

\begin{definitionbox}[Symétrie orthogonale]
Soit \(F\) un sous-espace de \(E\).  
La symétrie orthogonale par rapport à \(F\) est l'endomorphisme \(s\) défini par :
\[
x=x_F+x_{F^\perp}
\quad\Longrightarrow\quad
s(x)=x_F-x_{F^\perp},
\]
où :
\[
x_F\in F,
\qquad
x_{F^\perp}\in F^\perp.
\]
\end{definitionbox}

\begin{theorembox}
La symétrie orthogonale par rapport à \(F\) est une isométrie vectorielle.
\end{theorembox}

\begin{proof}
Soit :
\[
x=x_F+x_{F^\perp}.
\]
Alors :
\[
s(x)=x_F-x_{F^\perp}.
\]
Comme \(x_F\perp x_{F^\perp}\), on a :
\[
\norme{x}^2=\norme{x_F}^2+\norme{x_{F^\perp}}^2.
\]
De même :
\[
\norme{s(x)}^2
=
\norme{x_F-x_{F^\perp}}^2
=
\norme{x_F}^2+\norme{x_{F^\perp}}^2.
\]
Donc :
\[
\norme{s(x)}=\norme{x}.
\]
Ainsi \(s\) est une isométrie vectorielle.
\end{proof}

\begin{retenirbox}
Une réflexion est une symétrie orthogonale par rapport à un hyperplan.
\end{retenirbox}

%=========================================================
\section{Endomorphismes symétriques et théorème spectral}

\subsection{Endomorphismes symétriques}

\begin{definitionbox}[Endomorphisme symétrique]
Soit \(E\) un espace euclidien.  
Un endomorphisme \(u\in\Lcal(E)\) est dit symétrique si :
\[
\forall x,y\in E,\qquad
\scal{u(x)}{y}=\scal{x}{u(y)}.
\]
\end{definitionbox}

\begin{theorembox}[Caractérisation matricielle]
Dans une base orthonormée, \(u\) est symétrique si et seulement si sa matrice \(A\) vérifie :
\[
A^T=A.
\]
\end{theorembox}

\begin{proof}
Soit \(A\) la matrice de \(u\) dans une base orthonormée.  
Pour des colonnes de coordonnées \(X,Y\), on a :
\[
\scal{u(x)}{y}=(AX)^TY=X^TA^TY.
\]
Et :
\[
\scal{x}{u(y)}=X^T(AY)=X^TAY.
\]
L'égalité pour tous \(x,y\) équivaut à :
\[
X^TA^TY=X^TAY
\]
pour tous \(X,Y\), donc à :
\[
A^T=A.
\]
\end{proof}

\subsection{Orthogonalité des sous-espaces propres}

\begin{theorembox}
Si \(u\) est symétrique, alors les sous-espaces propres associés à des valeurs propres distinctes sont orthogonaux.
\end{theorembox}

\begin{proof}
Soient \(x\in E_\lambda(u)\) et \(y\in E_\mu(u)\), avec \(\lambda\neq\mu\).  
Alors :
\[
u(x)=\lambda x,
\qquad
u(y)=\mu y.
\]
Comme \(u\) est symétrique :
\[
\scal{u(x)}{y}=\scal{x}{u(y)}.
\]
Donc :
\[
\scal{\lambda x}{y}
=
\scal{x}{\mu y}.
\]
Par bilinéarité :
\[
\lambda\scal{x}{y}
=
\mu\scal{x}{y}.
\]
Donc :
\[
(\lambda-\mu)\scal{x}{y}=0.
\]
Comme \(\lambda\neq\mu\), on obtient :
\[
\scal{x}{y}=0.
\]
Donc :
\[
E_\lambda(u)\perp E_\mu(u).
\]
\end{proof}

\subsection{Théorème spectral}

\begin{theorembox}[Théorème spectral réel]
Tout endomorphisme symétrique d'un espace euclidien est diagonalisable dans une base orthonormée.

Sous forme matricielle : toute matrice symétrique réelle \(A\) est orthogonalement diagonalisable, c'est-à-dire qu'il existe une matrice orthogonale \(P\) et une matrice diagonale réelle \(D\) telles que :
\[
A=PDP^T.
\]
\end{theorembox}

\begin{proof}
On donne une démonstration adaptée au niveau prépa.

Soit \(u\) un endomorphisme symétrique de \(E\), avec \(\dim(E)=n\).  
On raisonne par récurrence sur \(n\).

Si \(n=1\), le résultat est immédiat.

Supposons \(n\geq2\).  
On admet ici le fait qu'un endomorphisme symétrique réel possède au moins une valeur propre réelle.  
Soit donc \(\lambda\) une valeur propre réelle, et soit \(e_1\) un vecteur propre unitaire associé :
\[
u(e_1)=\lambda e_1,
\qquad
\norme{e_1}=1.
\]

Considérons :
\[
F=e_1^\perp.
\]
Montrons que \(F\) est stable par \(u\).  
Soit \(x\in F\).  
On veut montrer que \(u(x)\in F\), c'est-à-dire :
\[
\scal{u(x)}{e_1}=0.
\]
Comme \(u\) est symétrique :
\[
\scal{u(x)}{e_1}
=
\scal{x}{u(e_1)}
=
\scal{x}{\lambda e_1}
=
\lambda\scal{x}{e_1}.
\]
Or \(x\in e_1^\perp\), donc :
\[
\scal{x}{e_1}=0.
\]
Ainsi :
\[
\scal{u(x)}{e_1}=0,
\]
donc \(u(x)\in F\).  
Le sous-espace \(F\) est stable par \(u\).

La restriction de \(u\) à \(F\) est encore symétrique.  
Par hypothèse de récurrence, il existe une base orthonormée de \(F\) formée de vecteurs propres de \(u_{\mid F}\).  
En ajoutant \(e_1\), on obtient une base orthonormée de \(E\) formée de vecteurs propres de \(u\).  
Donc \(u\) est diagonalisable dans une base orthonormée.
\end{proof}

\begin{retenirbox}
Une matrice symétrique réelle est toujours diagonalisable, et la matrice de passage peut être choisie orthogonale.
\[
A=A^T
\quad\Longrightarrow\quad
A=PDP^T.
\]
\end{retenirbox}

%=========================================================
\section{Endomorphismes symétriques positifs}

\begin{definitionbox}[Endomorphisme symétrique positif]
Un endomorphisme symétrique \(u\) est dit positif si :
\[
\forall x\in E,\qquad \scal{u(x)}{x}\geq0.
\]
Il est dit défini positif si :
\[
\forall x\neq0,\qquad \scal{u(x)}{x}>0.
\]
\end{definitionbox}

\begin{theorembox}
Soit \(u\) un endomorphisme symétrique.  
Alors \(u\) est positif si et seulement si toutes ses valeurs propres sont positives ou nulles.
\end{theorembox}

\begin{proof}
Par le théorème spectral, il existe une base orthonormée \((e_1,\dots,e_n)\) formée de vecteurs propres de \(u\).  
On note :
\[
u(e_i)=\lambda_i e_i.
\]

Soit :
\[
x=\sum_{i=1}^n x_ie_i.
\]
Alors :
\[
u(x)=\sum_{i=1}^n \lambda_i x_ie_i.
\]
Donc :
\[
\scal{u(x)}{x}
=
\sum_{i=1}^n \lambda_i x_i^2.
\]

Si tous les \(\lambda_i\) sont positifs ou nuls, alors :
\[
\scal{u(x)}{x}\geq0.
\]

Réciproquement, si \(u\) est positif, alors pour chaque vecteur propre unitaire \(e_i\),
\[
\lambda_i=\scal{u(e_i)}{e_i}\geq0.
\]
Donc toutes les valeurs propres sont positives ou nulles.
\end{proof}

\begin{corollaire}
Un endomorphisme symétrique est défini positif si et seulement si toutes ses valeurs propres sont strictement positives.
\end{corollaire}

\begin{proof}
Avec les notations précédentes :
\[
\scal{u(x)}{x}
=
\sum_{i=1}^n \lambda_i x_i^2.
\]
Cette quantité est strictement positive pour tout \(x\neq0\) si et seulement si tous les \(\lambda_i\) sont strictement positifs.
\end{proof}

%=========================================================
\section{Exercices progressifs avec corrigés}

\subsection*{Exercice 1 -- Vérifier un produit scalaire}

\begin{exercicebox}
Sur \(\R^2\), on définit :
\[
\scal{(x,y)}{(x',y')}=2xx'+xy'+x'y+3yy'.
\]
Montrer que cette application est un produit scalaire.
\end{exercicebox}

\begin{correctionbox}
La bilinéarité est immédiate car l'expression est linéaire par rapport à chaque couple de variables.

La symétrie est vérifiée :
\[
2xx'+xy'+x'y+3yy'
=
2x'x+x'y+xy'+3y'y.
\]

Il reste à vérifier la positivité définie.  
Pour \((x,y)\in\R^2\),
\[
\scal{(x,y)}{(x,y)}
=
2x^2+2xy+3y^2.
\]
On complète le carré :
\[
2x^2+2xy+3y^2
=
2\left(x+\frac y2\right)^2+\frac52y^2.
\]
Cette quantité est positive.  
Si elle est nulle, alors :
\[
2\left(x+\frac y2\right)^2=0
\quad\text{et}\quad
\frac52y^2=0.
\]
Donc \(y=0\), puis \(x=0\).  
Ainsi l'application est définie positive.  
C'est donc un produit scalaire.
\end{correctionbox}

\subsection*{Exercice 2 -- Orthogonal d'un plan}

\begin{exercicebox}
Dans \(\R^3\) muni du produit scalaire usuel, soit :
\[
F=\{(x,y,z)\in\R^3:\ x+y+z=0\}.
\]
Déterminer \(F^\perp\).
\end{exercicebox}

\begin{correctionbox}
Le plan \(F\) est donné par :
\[
x+y+z=0.
\]
Cette équation s'écrit :
\[
\scal{(x,y,z)}{(1,1,1)}=0.
\]
Donc :
\[
F=(1,1,1)^\perp.
\]
Ainsi :
\[
F^\perp=\Vect((1,1,1)).
\]
\end{correctionbox}

\subsection*{Exercice 3 -- Gram-Schmidt}

\begin{exercicebox}
Dans \(\R^3\), appliquer Gram-Schmidt à la famille :
\[
u_1=(1,1,0),\qquad u_2=(1,0,1),\qquad u_3=(0,1,1).
\]
\end{exercicebox}

\begin{correctionbox}
On pose :
\[
v_1=u_1=(1,1,0).
\]
Alors :
\[
\norme{v_1}^2=2.
\]

Ensuite :
\[
v_2=u_2-\frac{\scal{u_2}{v_1}}{\norme{v_1}^2}v_1.
\]
On calcule :
\[
\scal{u_2}{v_1}=1.
\]
Donc :
\[
v_2=(1,0,1)-\frac12(1,1,0)
=
\left(\frac12,-\frac12,1\right).
\]
On peut multiplier par \(2\) pour simplifier :
\[
w_2=(1,-1,2).
\]
Alors \(w_2\) est orthogonal à \(v_1\).

Pour le troisième vecteur, on utilise \(v_1\) et \(w_2\).  
On pose :
\[
v_3=u_3-\frac{\scal{u_3}{v_1}}{\norme{v_1}^2}v_1
-\frac{\scal{u_3}{w_2}}{\norme{w_2}^2}w_2.
\]
On calcule :
\[
\scal{u_3}{v_1}=1,
\qquad
\norme{v_1}^2=2.
\]
Et :
\[
\scal{u_3}{w_2}=(-1)+2=1,
\qquad
\norme{w_2}^2=1+1+4=6.
\]
Donc :
\[
v_3=(0,1,1)-\frac12(1,1,0)-\frac16(1,-1,2).
\]
On obtient :
\[
v_3=
\left(-\frac23,\frac23,\frac23\right).
\]
On peut multiplier par \(3/2\) :
\[
w_3=(-1,1,1).
\]
Une famille orthogonale associée est donc :
\[
(1,1,0),\qquad (1,-1,2),\qquad (-1,1,1).
\]
Après normalisation :
\[
e_1=\frac1{\sqrt2}(1,1,0),
\]
\[
e_2=\frac1{\sqrt6}(1,-1,2),
\]
\[
e_3=\frac1{\sqrt3}(-1,1,1).
\]
\end{correctionbox}

\subsection*{Exercice 4 -- Projection orthogonale}

\begin{exercicebox}
Dans \(\R^3\), soit :
\[
F=\Vect((1,1,0),(1,-1,2)).
\]
Calculer le projeté orthogonal de :
\[
x=(2,0,1)
\]
sur \(F\).
\end{exercicebox}

\begin{correctionbox}
On remarque que :
\[
u=(1,1,0),
\qquad
v=(1,-1,2)
\]
sont orthogonaux, car :
\[
\scal{u}{v}=1-1+0=0.
\]
Donc une base orthonormée de \(F\) est :
\[
e_1=\frac1{\sqrt2}(1,1,0),
\qquad
e_2=\frac1{\sqrt6}(1,-1,2).
\]
Alors :
\[
p_F(x)=\scal{x}{e_1}e_1+\scal{x}{e_2}e_2.
\]
On calcule :
\[
\scal{x}{e_1}
=
\frac1{\sqrt2}(2+0)=\sqrt2.
\]
Donc :
\[
\scal{x}{e_1}e_1=(1,1,0).
\]
Ensuite :
\[
\scal{x}{e_2}
=
\frac1{\sqrt6}(2+0+2)=\frac4{\sqrt6}.
\]
Donc :
\[
\scal{x}{e_2}e_2
=
\frac4{\sqrt6}\cdot\frac1{\sqrt6}(1,-1,2)
=
\frac23(1,-1,2).
\]
Ainsi :
\[
p_F(x)
=
(1,1,0)+\left(\frac23,-\frac23,\frac43\right)
=
\left(\frac53,\frac13,\frac43\right).
\]
\end{correctionbox}

\subsection*{Exercice 5 -- Distance à un plan}

\begin{exercicebox}
Dans \(\R^3\), calculer la distance du point :
\[
x=(1,2,3)
\]
au plan :
\[
F=\{(a,b,c):a+b+c=0\}.
\]
\end{exercicebox}

\begin{correctionbox}
On a :
\[
F=(1,1,1)^\perp.
\]
Donc \(F^\perp=\Vect((1,1,1))\).  
La distance de \(x\) à \(F\) est la norme de la projection de \(x\) sur \(F^\perp\).

La projection de \(x\) sur la droite \(\R n\), avec \(n=(1,1,1)\), vaut :
\[
p_{F^\perp}(x)=\frac{\scal{x}{n}}{\norme{n}^2}n.
\]
On calcule :
\[
\scal{x}{n}=1+2+3=6,
\qquad
\norme{n}^2=3.
\]
Donc :
\[
p_{F^\perp}(x)=2(1,1,1)=(2,2,2).
\]
La distance vaut :
\[
d(x,F)=\norme{p_{F^\perp}(x)}=\norme{(2,2,2)}=2\sqrt3.
\]
\end{correctionbox}

\subsection*{Exercice 6 -- Matrice orthogonale}

\begin{exercicebox}
Montrer que la matrice :
\[
A=
\frac{1}{\sqrt2}
\begin{pmatrix}
1&-1\\
1&1
\end{pmatrix}
\]
est orthogonale. Interpréter géométriquement.
\end{exercicebox}

\begin{correctionbox}
Les colonnes de \(A\) sont :
\[
C_1=\frac1{\sqrt2}\begin{pmatrix}1\\1\end{pmatrix},
\qquad
C_2=\frac1{\sqrt2}\begin{pmatrix}-1\\1\end{pmatrix}.
\]
On a :
\[
\norme{C_1}=1,
\qquad
\norme{C_2}=1,
\]
et :
\[
\scal{C_1}{C_2}
=
\frac12(-1+1)=0.
\]
Les colonnes forment donc une base orthonormée de \(\R^2\).  
La matrice est orthogonale.

De plus :
\[
\det(A)=1.
\]
C'est donc une rotation du plan.  
On reconnaît :
\[
A=
\begin{pmatrix}
\cos\theta&-\sin\theta\\
\sin\theta&\cos\theta
\end{pmatrix}
\]
avec :
\[
\cos\theta=\sin\theta=\frac1{\sqrt2}.
\]
Donc :
\[
\theta=\frac\pi4.
\]
C'est la rotation d'angle \(\pi/4\).
\end{correctionbox}

\subsection*{Exercice 7 -- Théorème spectral}

\begin{exercicebox}
Diagonaliser orthogonalement :
\[
A=
\begin{pmatrix}
2&1\\
1&2
\end{pmatrix}.
\]
\end{exercicebox}

\begin{correctionbox}
La matrice est symétrique réelle, donc orthogonalement diagonalisable.

On calcule :
\[
\chi_A(X)
=
\det
\begin{pmatrix}
X-2&-1\\
-1&X-2
\end{pmatrix}
=
(X-2)^2-1
=
X^2-4X+3.
\]
Donc :
\[
\chi_A(X)=(X-1)(X-3).
\]
Les valeurs propres sont \(1\) et \(3\).

Pour \(\lambda=1\),
\[
A-I=
\begin{pmatrix}
1&1\\
1&1
\end{pmatrix}.
\]
On obtient :
\[
x+y=0.
\]
Donc :
\[
E_1=\Vect((1,-1)).
\]
Un vecteur unitaire est :
\[
e_1=\frac1{\sqrt2}(1,-1).
\]

Pour \(\lambda=3\),
\[
A-3I=
\begin{pmatrix}
-1&1\\
1&-1
\end{pmatrix}.
\]
On obtient :
\[
x=y.
\]
Donc :
\[
E_3=\Vect((1,1)).
\]
Un vecteur unitaire est :
\[
e_2=\frac1{\sqrt2}(1,1).
\]

On pose :
\[
P=
\frac1{\sqrt2}
\begin{pmatrix}
1&1\\
-1&1
\end{pmatrix},
\qquad
D=
\begin{pmatrix}
1&0\\
0&3
\end{pmatrix}.
\]
Alors :
\[
A=PDP^T.
\]
\end{correctionbox}

%=========================================================
\newpage
\section{Grand devoir bilan final}

\begin{center}
\Large\textbf{Devoir bilan -- Espaces euclidiens}
\end{center}

\subsection*{Exercice 1 -- Produit scalaire et Gram-Schmidt}

Sur \(\R^3\), on considère le produit scalaire usuel et les vecteurs :
\[
u_1=(1,1,1),
\qquad
u_2=(1,1,0),
\qquad
u_3=(1,0,0).
\]

\begin{enumerate}
    \item Montrer que \((u_1,u_2,u_3)\) est une base de \(\R^3\).
    \item Appliquer le procédé de Gram-Schmidt à cette base.
    \item En déduire une base orthonormée de \(\R^3\).
\end{enumerate}

\subsection*{Exercice 2 -- Projection et distance}

On considère le plan :
\[
F=\Vect((1,1,0),(1,0,1)).
\]
Soit :
\[
x=(1,2,3).
\]

\begin{enumerate}
    \item Trouver une base orthonormée de \(F\).
    \item Calculer le projeté orthogonal de \(x\) sur \(F\).
    \item Calculer la distance de \(x\) à \(F\).
\end{enumerate}

\subsection*{Exercice 3 -- Matrice orthogonale}

Soit :
\[
A=
\begin{pmatrix}
0&-1&0\\
1&0&0\\
0&0&1
\end{pmatrix}.
\]

\begin{enumerate}
    \item Montrer que \(A\) est orthogonale.
    \item Calculer \(\det(A)\).
    \item Interpréter géométriquement l'endomorphisme associé.
\end{enumerate}

\subsection*{Exercice 4 -- Théorème spectral}

Soit :
\[
B=
\begin{pmatrix}
3&0&0\\
0&2&1\\
0&1&2
\end{pmatrix}.
\]

\begin{enumerate}
    \item Montrer que \(B\) est symétrique.
    \item Justifier que \(B\) est orthogonalement diagonalisable.
    \item Déterminer ses valeurs propres.
    \item Donner une base orthonormée de vecteurs propres.
    \item Écrire une diagonalisation orthogonale de \(B\).
\end{enumerate}

\subsection*{Exercice 5 -- Positivité}

Soit :
\[
C=
\begin{pmatrix}
2&1\\
1&2
\end{pmatrix}.
\]
\begin{enumerate}
    \item Montrer que \(C\) est symétrique.
    \item Déterminer ses valeurs propres.
    \item Montrer que \(C\) est définie positive.
    \item En déduire que :
    \[
    2x^2+2xy+2y^2>0
    \]
    pour tout \((x,y)\neq(0,0)\).
\end{enumerate}

\newpage

%=========================================================
\section{Correction détaillée du devoir bilan}

\subsection*{Correction de l'exercice 1}

On considère :
\[
u_1=(1,1,1),
\qquad
u_2=(1,1,0),
\qquad
u_3=(1,0,0).
\]

\textbf{1.}  
La matrice ayant ces vecteurs pour colonnes est :
\[
M=
\begin{pmatrix}
1&1&1\\
1&1&0\\
1&0&0
\end{pmatrix}.
\]
Son déterminant vaut :
\[
\det(M)
=
1\begin{vmatrix}1&0\\0&0\end{vmatrix}
-
1\begin{vmatrix}1&0\\1&0\end{vmatrix}
+
1\begin{vmatrix}1&1\\1&0\end{vmatrix}
=
0-0-1=-1.
\]
Il est non nul.  
Donc \((u_1,u_2,u_3)\) est une base de \(\R^3\).

\textbf{2. Gram-Schmidt.}

On pose :
\[
v_1=u_1=(1,1,1).
\]
Alors :
\[
\norme{v_1}^2=3.
\]

Ensuite :
\[
v_2=u_2-\frac{\scal{u_2}{v_1}}{\norme{v_1}^2}v_1.
\]
On a :
\[
\scal{u_2}{v_1}=2.
\]
Donc :
\[
v_2=(1,1,0)-\frac23(1,1,1)
=
\left(\frac13,\frac13,-\frac23\right).
\]
On peut prendre :
\[
w_2=(1,1,-2).
\]

Enfin :
\[
v_3=u_3-\frac{\scal{u_3}{v_1}}{\norme{v_1}^2}v_1
-\frac{\scal{u_3}{w_2}}{\norme{w_2}^2}w_2.
\]
On calcule :
\[
\scal{u_3}{v_1}=1,
\qquad
\norme{v_1}^2=3,
\]
\[
\scal{u_3}{w_2}=1,
\qquad
\norme{w_2}^2=1+1+4=6.
\]
Donc :
\[
v_3=(1,0,0)-\frac13(1,1,1)-\frac16(1,1,-2).
\]
On obtient :
\[
v_3=
\left(\frac12,-\frac12,0\right).
\]
On peut prendre :
\[
w_3=(1,-1,0).
\]

Une famille orthogonale est donc :
\[
(1,1,1),
\qquad
(1,1,-2),
\qquad
(1,-1,0).
\]

\textbf{3. Normalisation.}

On obtient :
\[
e_1=\frac1{\sqrt3}(1,1,1),
\]
\[
e_2=\frac1{\sqrt6}(1,1,-2),
\]
\[
e_3=\frac1{\sqrt2}(1,-1,0).
\]
La famille \((e_1,e_2,e_3)\) est une base orthonormée de \(\R^3\).

\subsection*{Correction de l'exercice 2}

On pose :
\[
a=(1,1,0),
\qquad
b=(1,0,1).
\]

On applique Gram-Schmidt.  
On prend :
\[
v_1=a=(1,1,0),
\qquad
\norme{v_1}^2=2.
\]
Puis :
\[
v_2=b-\frac{\scal{b}{v_1}}{\norme{v_1}^2}v_1.
\]
On a :
\[
\scal{b}{v_1}=1.
\]
Donc :
\[
v_2=(1,0,1)-\frac12(1,1,0)
=
\left(\frac12,-\frac12,1\right).
\]
On peut prendre :
\[
w_2=(1,-1,2).
\]
Alors :
\[
\norme{w_2}^2=1+1+4=6.
\]
Une base orthonormée de \(F\) est donc :
\[
e_1=\frac1{\sqrt2}(1,1,0),
\qquad
e_2=\frac1{\sqrt6}(1,-1,2).
\]

Soit :
\[
x=(1,2,3).
\]
Le projeté orthogonal est :
\[
p_F(x)=\scal{x}{e_1}e_1+\scal{x}{e_2}e_2.
\]
On calcule :
\[
\scal{x}{e_1}=\frac1{\sqrt2}(1+2)=\frac3{\sqrt2}.
\]
Donc :
\[
\scal{x}{e_1}e_1
=
\frac32(1,1,0)
=
\left(\frac32,\frac32,0\right).
\]
Ensuite :
\[
\scal{x}{e_2}
=
\frac1{\sqrt6}(1-2+6)
=
\frac5{\sqrt6}.
\]
Donc :
\[
\scal{x}{e_2}e_2
=
\frac56(1,-1,2)
=
\left(\frac56,-\frac56,\frac53\right).
\]
Ainsi :
\[
p_F(x)
=
\left(\frac32+\frac56,\frac32-\frac56,\frac53\right)
=
\left(\frac73,\frac23,\frac53\right).
\]

La distance est :
\[
d(x,F)=\norme{x-p_F(x)}.
\]
On calcule :
\[
x-p_F(x)
=
\left(1-\frac73,2-\frac23,3-\frac53\right)
=
\left(-\frac43,\frac43,\frac43\right).
\]
Donc :
\[
d(x,F)
=
\sqrt{
\frac{16}{9}+\frac{16}{9}+\frac{16}{9}
}
=
\frac{4\sqrt3}{3}.
\]

\subsection*{Correction de l'exercice 3}

Les colonnes de \(A\) sont :
\[
C_1=(0,1,0),
\qquad
C_2=(-1,0,0),
\qquad
C_3=(0,0,1).
\]
Elles sont unitaires et deux à deux orthogonales.  
Donc \(A\) est orthogonale.

On calcule :
\[
\det(A)=1.
\]
Cette matrice représente une rotation d'angle \(\pi/2\) dans le plan \((xOy)\), laissant fixe l'axe \(Oz\).

\subsection*{Correction de l'exercice 4}

La matrice :
\[
B=
\begin{pmatrix}
3&0&0\\
0&2&1\\
0&1&2
\end{pmatrix}
\]
est symétrique, car :
\[
B^T=B.
\]
Donc, par le théorème spectral, elle est orthogonalement diagonalisable.

Calculons son polynôme caractéristique :
\[
\chi_B(X)
=
\det
\begin{pmatrix}
X-3&0&0\\
0&X-2&-1\\
0&-1&X-2
\end{pmatrix}.
\]
Donc :
\[
\chi_B(X)
=
(X-3)\left((X-2)^2-1\right).
\]
Or :
\[
(X-2)^2-1=X^2-4X+3=(X-1)(X-3).
\]
Ainsi :
\[
\chi_B(X)=(X-3)^2(X-1).
\]
Les valeurs propres sont :
\[
1
\quad\text{et}\quad
3.
\]

Pour \(\lambda=1\),
\[
B-I=
\begin{pmatrix}
2&0&0\\
0&1&1\\
0&1&1
\end{pmatrix}.
\]
Le système donne :
\[
x=0,
\qquad
y+z=0.
\]
Donc :
\[
E_1=\Vect((0,1,-1)).
\]
Un vecteur unitaire est :
\[
e_1=\frac1{\sqrt2}(0,1,-1).
\]

Pour \(\lambda=3\),
\[
B-3I=
\begin{pmatrix}
0&0&0\\
0&-1&1\\
0&1&-1
\end{pmatrix}.
\]
Le système donne :
\[
-y+z=0,
\]
donc \(y=z\).  
Ainsi :
\[
E_3=\{(x,y,y):x,y\in\R\}
=
\Vect((1,0,0),(0,1,1)).
\]
Une base orthonormée de \(E_3\) est :
\[
e_2=(1,0,0),
\qquad
e_3=\frac1{\sqrt2}(0,1,1).
\]

On pose :
\[
P=
\begin{pmatrix}
0&1&0\\
\frac1{\sqrt2}&0&\frac1{\sqrt2}\\
-\frac1{\sqrt2}&0&\frac1{\sqrt2}
\end{pmatrix},
\qquad
D=
\begin{pmatrix}
1&0&0\\
0&3&0\\
0&0&3
\end{pmatrix}.
\]
Alors :
\[
B=PDP^T.
\]

\subsection*{Correction de l'exercice 5}

La matrice :
\[
C=
\begin{pmatrix}
2&1\\
1&2
\end{pmatrix}
\]
est symétrique.

On a déjà calculé dans le cours :
\[
\chi_C(X)=(X-1)(X-3).
\]
Les valeurs propres sont donc :
\[
1
\quad\text{et}\quad
3.
\]
Elles sont toutes strictement positives.  
Comme \(C\) est symétrique réelle, elle est définie positive.

Pour tout \((x,y)\neq(0,0)\),
\[
\begin{pmatrix}x&y\end{pmatrix}
C
\begin{pmatrix}x\\y\end{pmatrix}
=
2x^2+2xy+2y^2.
\]
Comme \(C\) est définie positive, cette quantité est strictement positive pour tout \((x,y)\neq(0,0)\).  
Donc :
\[
2x^2+2xy+2y^2>0.
\]

%=========================================================
\newpage
\section{Fiche de synthèse finale}

\subsection*{1. Produit scalaire}

Un produit scalaire sur un espace vectoriel réel \(E\) est une application :
\[
\scal{\cdot}{\cdot}:E\times E\to\R
\]
bilinéaire, symétrique, positive et définie positive.

La norme associée est :
\[
\norme{x}=\sqrt{\scal{x}{x}}.
\]

\subsection*{2. Cauchy-Schwarz}

\[
|\scal{x}{y}|\leq\norme{x}\norme{y}.
\]

Égalité si et seulement si \(x\) et \(y\) sont colinéaires.

\subsection*{3. Identités}

Polarisation :
\[
\scal{x}{y}
=
\frac12(\norme{x+y}^2-\norme{x}^2-\norme{y}^2).
\]

Autre forme :
\[
\scal{x}{y}
=
\frac14(\norme{x+y}^2-\norme{x-y}^2).
\]

Parallélogramme :
\[
\norme{x+y}^2+\norme{x-y}^2
=
2\norme{x}^2+2\norme{y}^2.
\]

\subsection*{4. Orthogonalité}

\[
x\perp y
\Longleftrightarrow
\scal{x}{y}=0.
\]

Pythagore :
\[
x\perp y
\Longrightarrow
\norme{x+y}^2=\norme{x}^2+\norme{y}^2.
\]

Orthogonal d'une partie :
\[
A^\perp=\{x\in E:\forall a\in A,\ \scal{x}{a}=0\}.
\]

\subsection*{5. Familles orthonormées}

Une famille orthonormée vérifie :
\[
\scal{e_i}{e_j}=\delta_{ij}.
\]

Toute famille orthogonale de vecteurs non nuls est libre.

Dans une base orthonormée :
\[
x=\sum_{k=1}^n\scal{x}{e_k}e_k.
\]

Parseval :
\[
\norme{x}^2=\sum_{k=1}^n\scal{x}{e_k}^2.
\]

\subsection*{6. Gram-Schmidt}

À partir d'une famille libre \((u_1,\dots,u_n)\), on construit :
\[
v_1=u_1,
\]
\[
v_k=u_k-\sum_{j=1}^{k-1}
\frac{\scal{u_k}{v_j}}{\norme{v_j}^2}v_j.
\]
Puis :
\[
e_k=\frac{v_k}{\norme{v_k}}.
\]

\subsection*{7. Sous-espaces orthogonaux}

Pour tout sous-espace \(F\) :
\[
E=F\oplus F^\perp.
\]

\[
\dim(F)+\dim(F^\perp)=\dim(E).
\]

\[
(F^\perp)^\perp=F.
\]

\subsection*{8. Projection orthogonale}

Si \((e_1,\dots,e_p)\) est une base orthonormée de \(F\), alors :
\[
p_F(x)=\sum_{k=1}^p\scal{x}{e_k}e_k.
\]

Sur une droite \(\R a\) :
\[
p_{\R a}(x)=\frac{\scal{x}{a}}{\norme{a}^2}a.
\]

Distance à un sous-espace :
\[
d(x,F)=\norme{x-p_F(x)}.
\]

\subsection*{9. Matrices de Gram}

Dans une base \(\mathcal B\),
\[
G_{\mathcal B}=(\scal{e_i}{e_j}).
\]
Si \(X,Y\) sont les coordonnées de \(x,y\), alors :
\[
\scal{x}{y}=X^TG_{\mathcal B}Y.
\]

Dans une base orthonormée :
\[
G_{\mathcal B}=I_n.
\]

\subsection*{10. Matrices orthogonales}

\[
A\in O(n)
\Longleftrightarrow
A^TA=I_n.
\]

Équivalent :
les colonnes de \(A\) forment une base orthonormée.

\[
A^{-1}=A^T.
\]

\[
\det(A)=\pm1.
\]

\[
SO(n)=\{A\in O(n):\det(A)=1\}.
\]

\subsection*{11. Isométries}

Un endomorphisme \(u\) est une isométrie si :
\[
\norme{u(x)}=\norme{x}.
\]

Équivalent :
\[
\scal{u(x)}{u(y)}=\scal{x}{y}.
\]

Dans une base orthonormée, sa matrice est orthogonale.

\subsection*{12. Endomorphismes symétriques}

\[
u\text{ symétrique}
\Longleftrightarrow
\forall x,y,\quad
\scal{u(x)}{y}=\scal{x}{u(y)}.
\]

Dans une base orthonormée :
\[
u\text{ symétrique}
\Longleftrightarrow
A^T=A.
\]

Théorème spectral :
\[
A=A^T
\Longrightarrow
\exists P\in O(n),\exists D\text{ diagonale},\quad
A=PDP^T.
\]

\subsection*{13. Positivité}

Un endomorphisme symétrique \(u\) est positif si :
\[
\forall x,\quad \scal{u(x)}{x}\geq0.
\]

Il est défini positif si :
\[
\forall x\neq0,\quad \scal{u(x)}{x}>0.
\]

Pour un endomorphisme symétrique :
\[
u\text{ positif}
\Longleftrightarrow
\text{toutes ses valeurs propres sont positives ou nulles}.
\]

\[
u\text{ défini positif}
\Longleftrightarrow
\text{toutes ses valeurs propres sont strictement positives}.
\]

%=========================================================
\section*{Conclusion pédagogique}
\addcontentsline{toc}{section}{Conclusion pédagogique}

Les espaces euclidiens donnent un cadre algébrique rigoureux à la géométrie.  
Le produit scalaire permet de définir l'orthogonalité, les angles, les distances à des sous-espaces, les projections, les symétries orthogonales et les isométries.

Le chapitre repose sur une idée centrale : les bases orthonormées simplifient tous les calculs.  
Dans une base orthonormée, les coordonnées, les produits scalaires, les projections et les matrices orthogonales prennent leur forme la plus simple.

Enfin, le théorème spectral relie profondément l'orthogonalité et la réduction : toute matrice symétrique réelle est diagonalisable dans une base orthonormée.  
C'est l'un des résultats majeurs de l'algèbre linéaire euclidienne.

\end{document}
``
::contentReference[oaicite:3]{index=3}
`